\section{领导力方法：组织如何匹配系统复杂度}

\subsection{60 直报与群体决策机制}

Jensen 透露他有大规模 direct reports，并且不做传统 one-on-one 作为主要沟通模式，而是通过多方同时在场的问题讨论驱动决策。这种机制和 ``extreme co-design'' 是配套的：复杂问题需要跨域同时反馈，串行传递会显著降低决策质量。

\begin{importantbox}{组织结构服务于产出结构}
Jensen 的方法论是：公司不是部门集合，而是 ``生产产品的机器''。如果产品本身是跨层耦合系统，组织结构也必须支持跨层协同，否则会在接口处丢失信息。
\end{importantbox}

\subsection{Listening to the whispers}

访谈中的 ``listening to the whispers'' 是另一条关键方法：在大信号出现前捕捉弱信号。弱信号来自客户 workload、研究社区、供应链反馈和工程异常。它们不会自动形成决策，需要领导层主动组织吸收机制。

\begin{knowledgebox}{弱信号到决策的转化链}
有效转化通常分三步：
\begin{itemize}
    \item 先把弱信号结构化（问题是否重复出现、是否跨客户出现）。
    \item 再做小范围验证（实验、灰度、仿真）。
    \item 最后把通过验证的信号写进路线图与资源配置。
\end{itemize}
这套机制能解释为什么路线图会快速迭代。
\end{knowledgebox}

\subsection{公开压力下的领导者自我约束}

Jensen 提到，公开环境下的错误会被放大，这反过来形成对决策者的约束。对工程组织而言，这种约束的价值在于减少 ``不可验证的自信''，推动团队回到可测量、可复盘的决策框架。

\begin{figure}[H]
\centering
\includegraphics[width=0.92\textwidth]{frame-004.jpg}
\caption{访谈尾段：从领导力延伸到组织执行与个人工作方式\protect\footnotemark}
\end{figure}
\footnotetext{视频画面时间区间：02:07:40--02:08:05。画面对应 Jensen 讨论 AI 使用方式与组织执行。}

\subsection{本章小结}

领导力在这场对话里的定义很务实：把复杂问题拆成可协同、可复盘、可持续迭代的组织机制。Jensen 的方法不是神秘个人天赋，而是一套可观察的组织工程实践。

